\section{Results}

\subsection*{The support of Stadlmann ($k=49$, claimed $H_1\le240$)}

\begin{theorem}[\lbl{Computed}; exact, two independent programs]\label{thm:240}
For the parameters of \cite{Sta26} in the table above:
\begin{enumerate}
\item for $P$ the cap-free optimiser in $V_{20}$ ($707$ coefficients),
\[49J_T(P1_T)/I_T(P1_T)=4.0120389035\ldots;\]
\item for $P^\ast$ the top eigenvector of $(49G_J,G_I)$ in the same space,
rounded to $400$ significant bits per coefficient,
\[49J_T(P^\ast1_T)/I_T(P^\ast1_T)=4.0124198475\ldots.\]
\end{enumerate}
\end{theorem}

In both cases the two evaluators agree on $I_T$ and $J_T$ as exact rationals.
Without caps the same space gives $4.016176$, so the caps cost $0.0041$ for the
truncated optimiser and $0.0038$ for the optimised one: optimising on $T$
regains less than a tenth of the loss. Every point of $T$ has all
coordinates at most $0.6$ (\lbl{Lean-checked}, \texttt{stadlmann\_coord\_le})
and $T\subset L\R_{49}$; the dual-weight theorem of \cite[\S4]{Kim26},
applied as in \cite[\S5]{Kim26} and depending only on these two facts, gives
$\sup_FkJ_T/I_T\le4.14$ over \emph{all} $F$ (\lbl{Computed}, interval
arithmetic, \texttt{dual\_eps\_cert.certify\_box}, $337$ pieces). So the
margin of the variational step of \cite{Sta26} lies in $(0.0124,\,0.14]$.

\begin{corollary}[conditional; the analytic input is \lbl{Claimed}]
The criterion of \cite[Prop.~1]{Sta26} holds for $k=49$ with Stadlmann's
support, caps and parameters. With the \lbl{Lean-checked} admissible
$49$-tuple of diameter $240$, $H_1\le240$ follows from \cite[Prop.~1]{Sta26}
and his equidistribution theorem, which we have not verified.
\end{corollary}

\subsection*{The support of Kintali and of Song--Yue ($k=48$, claimed $H_1\le236$)}
Both preprints use the same support (the table above, called ``D182'' in
\cite{Kin26}) and, by their descriptions, the same public witness of degree
$25$, originally constructed for $k=50$ and evaluated at $k=48$ without
re-optimisation. Song and Yue print it as $1295$ integer triples
$(a,\alpha,c)$, with
\[
P(z)=\sum c\,\left(\tfrac{26}{25}-\textstyle\sum_iz_i\right)^a m_\alpha(z),\qquad
P_{\mathrm{SY}}(t)=P\!\left(\tfrac{26}{25U_{\mathrm{SY}}}t\right),\quad U_{\mathrm{SY}}=\tfrac{521}{2000},
\]
$m_\alpha$ the monomial symmetric polynomials, and they print the reference
integers $I_{\mathrm{ref}}=(98!)^2I_0^\ast$ and
$J_{\mathrm{ref}}=(99!)^2(26!)^2J^\ast$. Here $I_0$ is the \emph{cap-free} mass
on $\{\sum_it_i<U_{\mathrm{SY}}\}$, $J$ the capped numerator, both in their
$t$-coordinates, and the stars denote integer
coordinates $x=10^4t$. Their criterion is $48J/I_0>2001/2000$ (threshold $1$
in their normalisation), which implies
the capped criterion because $I_T\le I_0$ (\lbl{Literature},
\cite[\S4.12, App.~A, App.~C]{SY26}). They state that a fresh evaluation
reproduces $I_{\mathrm{ref}}$, while fresh regeneration of $J_{\mathrm{ref}}$
from the table ``remains incomplete''.

We extracted the table from the PDF using font sizes to read the exponents in
the partitions. All $1295$ rows are present, every $\alpha$ is nondecreasing
with even parts, and $\max(a+|\alpha|)=25$, as stated. Since all parts are
even, $P_{\mathrm{SY}}$, rewritten in $u=4t$, is exactly a polynomial in $L-P_1$ and the even power sums
$P_2,\dots,P_{24}$ ($1311$ terms); the conversion was checked exactly at random
rational points.

\begin{theorem}[\lbl{Computed}, exact rational arithmetic; items (1) and (2) each by two independent programs]\label{thm:sy}
For the witness $P_{\mathrm{SY}}$ of \cite{SY26} and the support above, with
$I_0$ and $J_T$ computed in the $t$-coordinates of \cite{SY26}:
\begin{enumerate}
\item $(98!)^2\,10^{192}\,5^{150}\,521^{50}\,I_0=I_{\mathrm{ref}}$ exactly
(an integer with $655$ digits);
\item $(99!)^2(26!)^2\,10^{196}\,5^{150}\,521^{50}\,J_T=J_{\mathrm{ref}}$
exactly (an integer with $715$ digits), where $J_T$ is the capped numerator;
\item hence, in the normalisation of \cite{SY26} (threshold $1$),
$48J_T/I_0=1.0006649215907\ldots>2001/2000$. In the Polymath scaling
(threshold $4$) the same ratio is $4.00265968636\ldots$; with the capped
denominator $48J_T/I_T=4.0030988047\ldots$, and without caps
$48J_0/I_0=4.0036348167\ldots$.
\end{enumerate}
\end{theorem}

The factors $10^{192}5^{150}521^{50}$ and $10^{196}5^{150}521^{50}$ come from the
change to integer coordinates when the polynomial is scaled by $65125^{25}$
(our reading of \cite[Eq.~(71)]{SY26}; $65125=25\cdot2605$), so items (1) and
(2) are equalities of integers, not numerical agreement. Item (1) was computed
by \texttt{songyue\_I0.py}, written for this purpose, which works in the monomial basis of
the table, integrating each product $m_\alpha m_\beta$ by a sum over partial
matchings of parts, and does not use the power-sum conversion;
\texttt{songyue\_J0.py} computes $J_0$ the same way. Item (2) was
computed by \texttt{capped\_check.py} on the converted polynomial ($56$ minutes
on four cores). Kintali \cite{Kin26} prints, in the normalisation with
threshold $1$, $48J_T/I_T=1.000774701187468\ldots$ and
$48J_0/I_0=1.000908704175987\ldots$, and the ratios $I_T/I_0=0.99989030538\ldots$
and $J_T/J_0=0.99975643873\ldots$; ours agree with all four to $12$
significant digits; for example
\[4\times1.000774701187468=4.003098804749\ldots,\]
and our capped ratio is $4.0030988047499\ldots$.

\emph{Second program (version 0.3).} The second evaluator,
\texttt{capped\_exact.py} (cap-free value minus $126$ inclusion--exclusion
correction pieces, each an exact rational; about $32$ core-hours), gives the
same $J_T$ as \texttt{capped\_check.py} as an exact rational, so item (2)
holds by two independent programs; its cap-free values equal $I_0$ and $J_0$
of the monomial-basis programs \texttt{songyue\_I0.py} and
\texttt{songyue\_J0.py} exactly, so item (1) does too (see
\texttt{code/}\allowbreak\texttt{songyue\_compare2.py} and
\texttt{code/}\allowbreak\texttt{cap6/}\allowbreak\texttt{compare2\_sy236.json}).
The capped denominator $I_T$, which enters only the ratio $48J_T/I_T$ and not
the criterion of \cite{SY26}, rests on \texttt{capped\_check.py} alone.
Proposition~\ref{prop:kin-own} below is also confirmed by both programs
(version 0.4).

\begin{proposition}[\lbl{Computed}, exact rational arithmetic, two independent programs]\label{prop:kin-own}
For our own cap-free optimiser in $V_{22}$ for this $(L,\tau)$ at $k=48$
($1041$ coefficients, rounded to $400$ significant bits),
$48J_T(P1_T)/I_T(P1_T)=4.0060814208\ldots>4$.
\end{proposition}

\texttt{capped\_check.py} and \texttt{capped\_exact.py} ($176$ correction
pieces, $15.5$ core-hours) give the same $I_T$ and $J_T$ as exact rationals
(\texttt{code/}\allowbreak\texttt{cap6/}\allowbreak\texttt{compare\_kin48\_D22r.json}).
Here the caps cost only $0.00033$ (the cap-free value is $4.0064115$). That our
value exceeds the witness's $4.0031$ is expected, since the witness was
optimised for $k=50$ and transported, while ours is optimised for $k=48$.

\begin{corollary}[conditional; the analytic input is \lbl{Claimed}]
For the published witness, the variational criteria of both \cite{Kin26}
($48J_T/I_T>4$) and \cite{SY26} ($48J_T/I_0>2001/2000$) hold at $k=48$. For
our independent test function only the capped criterion $48J_T/I_T>4$ was
computed, and it holds (two programs). With the
\lbl{Lean-checked} admissible $48$-tuple of diameter $236$, $H_1\le236$
follows from their equidistribution input, which we have not verified.
\end{corollary}

\subsection*{The draft of Charton et al.\ ($k=45$, claimed $H_1\le212$)}

\begin{proposition}[\lbl{Computed}; exact, two independent programs]\label{prop:212}
With the caps of \cite[Table~3]{Cha26} and $P$ our cap-free optimiser in
$V_{21}$ ($846$ coefficients, the dimension of the space
$\mathcal B_{21}$ of \cite[Thm.~11.1]{Cha26}),
$45J_T(P1_T)/I_T(P1_T)=4.0027351128\ldots>4$.
\end{proposition}

The authors report $4.00438$ with their own optimised function (\lbl{Claimed});
our value is lower, as expected for a function optimised without caps, and
already above $4$. The draft, however, also states caps in its Lemma~9.1,
namely $C_1=C_2=31/50$ and $C_m=17/25$ for $m\ge3$, which are smaller than
those of Table~3 for every $m\ge1$. With these the conclusion changes.

\begin{proposition}[\lbl{Computed}, ball arithmetic]\label{prop:ceiling}
Let $T'$ be the support with the caps of Lemma~9.1 of \cite{Cha26} and
$V=V_{21}$. Then $3.99\,G_I-45\,G_J$ is positive definite on
$V$, so
\[
\frac{45J_{T'}(P1_{T'})}{I_{T'}(P1_{T'})}<3.99\qquad\text{for every nonzero }P\in V.
\]
\end{proposition}

The weaker statement with $4$ in place of $3.99$ was certified twice, at
$5000$ bits (blocks halved) and at $6500$ bits (blocks split $1:2$); the
statement with $3.99$ was certified once, at $6500$ bits. Both runs use the
same saved weights, which agree exactly with \texttt{capped\_check.py} at a
random integer point. The best $P$ in $V$ gives
about $3.98555$ (floating point, from the Gram matrices). So within $V$, a
space of the size the authors use, the caps of Lemma~9.1 cannot give their
$4.00438$, which is consistent with the caps of Table~3. We report this as a
located inconsistency for the authors to resolve, not as an error in their
result. We have not checked that their space coincides with ours beyond its
dimension, and \texttt{capped\_exact.py} has not yet re-evaluated the random
point used for the weights.

